\section{Derived logical gates}

Derived gates are QPU instructions implemented as reversible target
mutations. Inputs are preserved. If target $t$ starts in zero, AND/OR/XOR
write the familiar classical function. If it does not, the true behavior is
the XOR embedding shown below.

\begin{beginnerbox}[One shape for every two-input gate]
AND, NAND, OR, and XOR all follow $t' = t\oplus f(a,b)$. The two inputs $a$
and $b$ (\textbf{Control A} and \textbf{Control B} in the workbench) are read
and left unchanged. The target $t$ (\textbf{Target particle}) is the output
workspace and is the only wire that changes. Only $f$ differs between the
four gates. To get the ordinary classical answer, start the target at 0 with
\code{SET Out 0p}.
\end{beginnerbox}

\GateHeading{Logical NOT}{NOT}{\GateSegment{$\neg$}}

\[
t'\;=\;t\oplus1.
\]
NOT is the derived logical spelling of a bit flip. In the current simulator it
uses the same X matrix on the output target. The required input establishes the
one-input instruction shape, but a distinct input is not copied into the
output; use the same token on both sides.
\begin{lstlisting}
NOT -I Target -O Target
\end{lstlisting}
\begin{center}
\begin{tabular}{c|c}\toprule $t$&$t'$\\\midrule0&1\\1&0\\\bottomrule\end{tabular}
\end{center}
The current gate implementation mutates the output target. Use the same token
for input and output; a distinct input token is not copied.

\GateHeading{Reversible AND}{AND}{\DoubleControlledSegment}

\[
t'=t\oplus(a\land b).
\]
AND flips the target exactly on the basis components where both controls are
one. This is Toffoli-style behavior. It remains linear on superpositions and
may entangle the target with its controls. Calling it AND is convenient only
when the target is deliberately initialized.
\begin{lstlisting}
SET Out 0p
AND -I A B -O Out
\end{lstlisting}
\begin{center}
\begin{tabular}{cc|c|c}\toprule $a$&$b$&$t=0$&$t'=a\land b$\\\midrule
0&0&0&0\\0&1&0&0\\1&0&0&0\\1&1&0&1\\\bottomrule
\end{tabular}
\end{center}

\GateHeading{Reversible NAND}{NAND}{\DoubleControlledSegment}

\[
t'=t\oplus(a\land b)\oplus1.
\]
NAND first performs the all-controls-active mutation and then inverts the
target. With target zero it produces the familiar NAND column. With target one
it instead produces AND, illustrating why reversible target state must always
be included in the full mathematical definition.
\begin{lstlisting}
SET Out 0p
NAND -I A B -O Out
\end{lstlisting}
\begin{center}
\begin{tabular}{cc|c|c}\toprule $a$&$b$&$t=0$&$t'=\neg(a\land b)$\\\midrule
0&0&0&1\\0&1&0&1\\1&0&0&1\\1&1&0&0\\\bottomrule
\end{tabular}
\end{center}

\GateHeading{Reversible OR}{OR}{\DoubleControlledSegment}

\[
t'=t\oplus(a\lor b).
\]
OR uses an ``any control is one'' predicate. If either or both controls are
active, the target is flipped. This direct derived operation is equivalent on
basis states to combining XOR and AND through
$a\lor b=a\oplus b\oplus(a\land b)$.
\begin{lstlisting}
SET Out 0p
OR -I A B -O Out
\end{lstlisting}
\begin{center}
\begin{tabular}{cc|c|c}\toprule $a$&$b$&$t=0$&$t'=a\lor b$\\\midrule
0&0&0&0\\0&1&0&1\\1&0&0&1\\1&1&0&1\\\bottomrule
\end{tabular}
\end{center}
For bits, $a\lor b=a\oplus b\oplus(a\land b)$.

\GateHeading{Reversible XOR}{XOR}{\DoubleControlledSegment}

\[
t'=t\oplus a\oplus b.
\]
XOR flips the target when an odd number of controls is one. With two inputs,
this is the ordinary exclusive-or predicate. The parity definition explains
why the 11 row does not flip: two active controls have even parity.
\begin{lstlisting}
SET Out 0p
XOR -I A B -O Out
\end{lstlisting}
\begin{center}
\begin{tabular}{cc|c|c}\toprule $a$&$b$&$t=0$&$t'=a\oplus b$\\\midrule
0&0&0&0\\0&1&0&1\\1&0&0&1\\1&1&0&0\\\bottomrule
\end{tabular}
\end{center}

